% answers_short.tex —— 2015 年·安徽卷·理科 · 简短答案速查
% 由 tools/make_short_answers.py 自 answers.tex 的 \gksolhead/\gkanswer 提取，
% 并与 2015/2015_数学真题_安徽理_答案要点.md 逐题比对后生成，勿手改（改源 MD 后重跑）。

\begin{gknotice}
\textbf{使用说明}：本节只给最终答案便于快速核对，完整推导见解析版。答案为编写者独立推导，\textbf{非官方标准答案}。
\end{gknotice}

\gkpart{选择题}
\begin{center}\begin{tabular}{ccccccccccc}
\toprule
题号 & 1 & 2 & 3 & 4 & 5 & 6 & 7 & 8 & 9 & 10 \\
\midrule
答案 & B & A & A & C & D & C & B & D & C & A \\
\bottomrule
\end{tabular}\end{center}

\gkpart{填空题}
\begin{itemize}[leftmargin=2.6em]
  \item[\textbf{11.}] \(35\)
  \item[\textbf{12.}] \(6\)
  \item[\textbf{13.}] \(4\)
  \item[\textbf{14.}] \(S_n=2^n-1\)
  \item[\textbf{15.}] ①③④⑤。
\end{itemize}

\gkpart{解答题}
\begin{itemize}[leftmargin=2.6em]
  \item[\textbf{16.}] \(AD=\sqrt{10}\)。
  \item[\textbf{17.}] (1) \(\dfrac{3}{10}\)；(2) \(P(X=200)=\dfrac{1}{10}\)、\(P(X=300)=\dfrac{3}{10}\)、\(P(X=400)=\dfrac{3}{5}\)，均值 \(E(X)=350\)（元）
  \item[\textbf{18.}] (1) \(x_n=\dfrac{n}{n+1}\)；(2) 成立（\(T_n=\prod_{k=1}^{n}\left(\dfrac{2k-1}{2k}\right)^2\)；对 \(k\geqslant 2\) 逐项用 \(\left(\dfrac{2k-1}{2k}\right)^2>\dfrac{k-1}{k}\)（分子比右端恰多 1）连乘作望远镜放缩，得 \(T_n>\dfrac{1}{4n}\)；\(n=1\) 时 \(T_1=\dfrac{1}{4}=\dfrac{1}{4\times 1}\) 取等）
  \item[\textbf{19.}] (1) 成立（\(F\) 为 \(CD_1\) 中点，\(EF\) 为 \(\triangle CD_1B_1\) 的中位线，故 \(EF\parallel B_1C\)）；(2) \(\dfrac{\sqrt{6}}{3}\)
  \item[\textbf{20.}] (1) \(e=\dfrac{2\sqrt{5}}{5}\)；(2) \(\dfrac{x^2}{45}+\dfrac{y^2}{9}=1\)
  \item[\textbf{21.}] (1) 当 \(a\leqslant -2\) 时 \(f(\sin x)\) 在 \(\left(-\dfrac{\pi}{2},\dfrac{\pi}{2}\right)\) 内单调递增、无极值；当 \(a\geqslant 2\) 时单调递减、无极值；当 \(-2<a<2\) 时先减后增，有极小值 \(b-\dfrac{a^2}{4}\)；(2) \(D=\max\left\{|(a+b)-(a_0+b_0)|,\ |(b-a)-(b_0-a_0)|\right\}\)；(3) 最大值 \(1\)
\end{itemize}
